\chapter{The Python Paradox}


            
              
                

August 2004

In a recent
                  \href{http://www.paulgraham.com/gh.html}{talk} I said something that upset a lot
                  of people: that you could get smarter programmers to work on a
                  Python project than you could to work on a Java project.

I
                  didn't mean by this that Java programmers are dumb. I meant
                  that Python programmers are smart. It's a lot of work to learn
                  a new programming language. And people don't learn Python
                  because it will get them a job; they learn it because they
                  genuinely like to program and aren't satisfied with the
                  languages they already know.

Which makes them
                  exactly the kind of programmers companies should want to hire.
                  Hence what, for lack of a better name, I'll call the Python
                  paradox: if a company chooses to write its software in a
                  comparatively esoteric language, they'll be able to hire
                  better programmers, because they'll attract only those who
                  cared enough to learn it. And for programmers the paradox is
                  even more pronounced: the language to learn, if you want to
                  get a good job, is a language that people don't learn merely
                  to get a job.

Only a few companies have been smart
                  enough to realize this so far. But there is a kind of
                  selection going on here too: they're exactly the companies
                  programmers would most like to work for. Google, for example.
                  When they advertise Java programming jobs, they also want
                  Python experience.

A friend of mine who knows
                  nearly all the widely used languages uses Python for most of
                  his projects. He says the main reason is that he likes the way
                  source code looks. That may seem a frivolous reason to choose
                  one language over another. But it is not so frivolous as it
                  sounds: when you program, you spend more time reading code
                  than writing it. You push blobs of source code around the way
                  a sculptor does blobs of clay. So a language that makes source
                  code ugly is maddening to an exacting programmer, as clay full
                  of lumps would be to a sculptor.

At the mention of
                  ugly source code, people will of course think of Perl. But the
                  superficial ugliness of Perl is not the sort I mean. Real
                  ugliness is not harsh-looking syntax, but having to build
                  programs out of the wrong concepts. Perl may look like a
                  cartoon character swearing, but there are
                  \href{http://www.paulgraham.com/icad.html}{cases} where it surpasses Python
                  conceptually.

So far, anyway. Both languages are of
                  course \href{http://www.paulgraham.com/hundred.html}{moving} targets. But they
                  share, along with Ruby (and Icon, and Joy, and J, and Lisp,
                  and Smalltalk) the fact that they're created by, and used by,
                  people who really care about programming. And those tend to be
                  the ones who do it well.


              
            
          

            
              
            
            
              
                \href{http://www.fazlamesai.net/modules.php?file=article\&name=News\&sid=2330}{Turkish Translation}

              
              
              
                \href{http://www.shiro.dreamhost.com/scheme/trans/pypar-j.html}{Japanese Translation}

              
            
            
              
            
            
              
            
            
              
                \href{http://www.sounerd.com.br/index.php?option=com\_content\&task=view\&id=191\&Itemid=43}{Portuguese Translation}

              
              
              
                \href{http://www.invece.org/translations/pparadox.html}{Italian Translation}

              
            
            
              
            
            
              
            
            
              
                \href{http://www.pdembinski.konin.lm.pl/python\_paradox.html}{Polish Translation}

              
              
              
                \href{http://ro.goobix.com/pg/pypar/}{Romanian Translation}

              
            
            
              
            
            
              
            
            
              
                \href{http://m0sia.ru/graham/pythonparadox}{Russian Translation}

              
              
              
                \href{http://www.fduran.com/wordpress/?p=23}{Spanish Translation}

              
            
            
              
            
            
              
            
            
              
                \href{http://w2.syronex.com/jmr/python-paradox}{French Translation}

              
              
              
                \href{http://www.avilpage.com/2014/12/python-paradox.html}{Telugu Translation}

              
            
            
              
            
          
